\scanpage{12}
\mspage{2}

est couvrant.

\uline{Hyp 2}\unsure{a} Pour $\alpha,\beta\in\mathrm{Ob}\,A$, l'ens des \deleted{flèches} diagrammes
\diagram{\begin{tikzcd}[ampersand replacement=\&,column sep=small,row sep=small]\&\alpha\\\gamma\arrow[ur]\arrow[dr]\&\\\&\beta\end{tikzcd}}
dans $A$ (i.e. l'ens. des objets de $A_{/\alpha\times\beta}$, où $\alpha\times\beta$ est le produit cartésien dans $\widehat A$), indexant les flèches correspondantes $U_\gamma\longrightarrow U_\alpha\times U_\beta$, donne une famille couvrant $U_\alpha\times U_\beta$.

[Idem pour une suite finie \inserted{(non vide)} $\alpha_0,\alpha_1,\ldots,\alpha_n$ ($n\in\mathbb N^{\deleted{\unsure{*}}}$) d'objets de $A$, et les \deleted{diagr.} flèches
\[U_\gamma\longrightarrow U_{\alpha_0}\times U_{\alpha_1}\times\cdots\times U_{\alpha_n}\]
indexées par les objets de $A_{/\alpha_0\times\alpha_1\times\cdots\alpha_n}$.]
\margin{[Hyp. 2 généralisé]}

Pour $n=0$, on retrouve Hyp 1', pour $n=1$, la formulation initiale de Hyp 2. Mais Hyp 2 pour $n=1$ implique-t-il Hyp 2 pour \deleted{$\unsure{n>1}$ \unsure{quelconque}} $n\geq1$. On voit qu'il en est ainsi, par récurrence, en utilisant le fait que la famille composée de deux familles couvrantes
\diagram{\begin{tikzcd}[ampersand replacement=\&,column sep=large]
U_{\gamma'}\arrow[r,"\varphi(v)"]\&U_\gamma\times U_{\alpha_{n+1}}\arrow[r,"\varphi(u)\times\mathrm{id}_{U_{\alpha_{n+1}}}"]\&(U_{\alpha_0}\times\cdots\times U_{\alpha_n})\times U_{\alpha_{n+1}}
\end{tikzcd}}
($v$ parcourant les flèches $\gamma'\longrightarrow\gamma\times\alpha_{n+1}$, et $u$ les flèches $\gamma\longrightarrow\alpha_0\times\cdots\times\alpha_n$) est couvrante. Ainsi Hyp 1' \deleted{implique} \inserted{$+$} Hyp 2\unsure{a}
\margin{11}
